\section*{Capítulo \ref{ch:b1:rifts-ridges-margins} --- Fosas, dorsales y márgenes pasivos}

\begin{solution}{exo:b1:rifts-ridges-margins:1}
(a) El dique cubriría en planta $\qty{60000}{m}\times\qty{8}{m} = \qty{4.8e5}{m^2}$, así que su altura media sería
$\qty{2.5e9}{m^3}/\qty{4.8e5}{m^2} =
\qty{5200}{m}$, unos \qty{5}{km}: el dique atravesó buena parte de la \omterm{def:b1:layered-earth:layers}{corteza}.
(b) $\qty{8}{m}/\qty{0.018}{m.yr^{-1}} = 444$ años, unos cuatro siglos.
\end{solution}

\begin{solution}{exo:b1:rifts-ridges-margins:2}
(a) El \omterm{def:b1:cross-sections:throw}{salto horizontal} es el salto dividido por la tangente del \omterm{def:b1:geological-maps:attitude}{buzamiento}:
$3.0/\tan\ang{60} = \qty{1.73}{km}$ y $0.8/\tan\ang{70} = \qty{0.29}{km}$;
la extensión de la sección es \qty{2.02}{km}. (b) Antes del fallamiento la
sección medía $25 - 2.02 = \qty{22.98}{km}$, unos \qty{23}{km}: se ha alargado en
unos \qty{9}{\percent}.
\end{solution}

\begin{solution}{exo:b1:rifts-ridges-margins:3}
Ultralenta (por debajo de \qty{20}{mm.yr^{-1}}): la dorsal de Gakkel (12), el mar Rojo central
(15) y la dorsal de Reykjanes (19), esta última cerca del límite. Lenta (de 20 a
\qty{55}{mm.yr^{-1}}): la dorsal de Carlsberg occidental (23) y el Atlántico Sur (33).
Rápida (por encima de \qty{75}{mm.yr^{-1}}): la dorsal del Pacífico oriental (143). Las
\omterm{def:b1:rifts-ridges-margins:ridge-classes}{dorsales ultralentas} y lentas deben tener un valle axial entre fallas normales; la
dorsal del Pacífico oriental, un alto axial.
\end{solution}

\begin{solution}{exo:b1:rifts-ridges-margins:4}
(a) $\qty{3.0}{km^2}\times\qty{7.1}{km} = \qty{21.3}{km^3}$, unos
\qty{21}{km^3} al año. (b) $\qty{3.028e8}{km^2}/\qty{3.0}{km^2.yr^{-1}} =
\qty{1.01e8}{yr}$, \qty{101}{Myr}. El fondo más antiguo, de unos \qty{190}{Myr}, es
más antiguo porque el fondo no se destruye en orden de edad: la subducción consume
fondo de todas las edades, parte del fondo antiguo sobrevive donde ninguna fosa lo ha alcanzado,
y la velocidad de producción ha variado.
\end{solution}

\begin{solution}{exo:b1:rifts-ridges-margins:5}
(a) $2500 + 350\times2 = \qty{3200}{m}$, $2500 + 350\times6 = \qty{4600}{m}$,
$2500 + 350\times8 = \qty{5300}{m}$. (b) $t = \bigl((d - 2500)/350\bigr)^2$:
$(2100/350)^2 = \qty{36}{Myr}$ y $(2800/350)^2 = \qty{64}{Myr}$. (c) El
fondo se hunde \qty{1400}{m} entre 4 y \qty{36}{Myr}, pero solo \qty{700}{m}
entre 36 y \qty{64}{Myr}, en un lapso de tiempo parecido.
\end{solution}

\begin{solution}{exo:b1:rifts-ridges-margins:6}
El coeficiente de dilatación efectivo es
$\num{3.2e-5}\times3300/2300 = \qty{4.591e-5}{K^{-1}}$, y
$\sqrt{\kappa\times\qty{1}{Myr}} = \sqrt{\num{e-6}\times\num{3.156e13}} =
\qty{5618}{m}$. Entonces $C = (2/\sqrt\pi)\times\num{4.591e-5}\times1350\times5618 =
1.1284\times\num{4.591e-5}\times1350\times5618 = \qty{393}{m.Myr^{-1/2}}$. Para un \omterm{def:b1:layered-earth:layers}{manto}
\qty{250}{K} más frío, $C$ escala con $T_m - T_0$:
$393\times1100/1350 = \qty{320}{m.Myr^{-1/2}}$. El fondo de un océano más frío
se hunde más despacio y sigue la constante menor. El \num{350} ajustado del
capítulo no es un caso así: queda dentro de la incertidumbre de los parámetros.
\end{solution}

\begin{solution}{exo:b1:rifts-ridges-margins:7}
(a) $\qty{e-6}{cal.cm^{-2}.s^{-1}} = \num{4.184e-6}\times\num{e4} =
\qty{4.184e-2}{W.m^{-2}} = \qty{41.84}{mW.m^{-2}}$, así que la constante es
$11.3\times41.84 = \qty{473}{mW.m^{-2}.Myr^{1/2}}$, un \qty{3.5}{\percent} por debajo de
los 490 de \cref{ch:b1:heat-budget}: los dos ajustes concuerdan. (b) $q = 473$,
$473/\sqrt{10} = 150$ y $473/\sqrt{150} = \qty{38.6}{mW.m^{-2}}$. La
ley de la raíz cuadrada sigue bajando, mientras que el modelo de placa se nivela en una constante:
a \qty{150}{Myr} el valor de la raíz cuadrada es demasiado bajo, como la ley de
profundidad se va demasiado hacia abajo.
\end{solution}

\begin{solution}{exo:b1:rifts-ridges-margins:8}
Tómese el segmento norte en $x = 0$ y el sur en $x = -40$
(km), la transformante sobre $y = 0$ entre ellos. (a) Justo al norte de la transformante el
fondo yace al oeste del segmento norte y se mueve hacia el oeste; justo al sur de ella, yace
al este del segmento sur y se mueve hacia el este. Visto desde cualquiera de los dos lados, el
lado lejano se mueve a la izquierda: el deslizamiento es levógiro (sinistral). (b) La
dorsal parece haber sido desplazada dextrógiramente, con su parte sur al
oeste: el sentido opuesto, como predice \cref{prop:b1:rifts-ridges-margins:transform}.
(c) Los terremotos ocurren solo en los \qty{40}{km} entre los dos extremos de la dorsal,
donde los dos lados pertenecen a placas distintas; más allá, a lo largo de las
\omterm{def:b1:rifts-ridges-margins:transform}{zonas de fractura}, ambos lados pertenecen a la misma placa y se mueven juntos.
\end{solution}

\begin{solution}{exo:b1:rifts-ridges-margins:9}
La velocidad media es \qty{10.6}{mm.yr^{-1}}, \qty{10.6}{km.Myr^{-1}}. (a) El escalón de edad es
$350/10.6 = \qty{33.0}{Myr}$. (b) Al final de la transformante, el fondo
de \qty{33.0}{Myr} ($2500 + 350\sqrt{33.0} = \qty{4511}{m}$) se enfrenta a fondo de
edad 0 (\qty{2500}{m}): un escalón de unos \qty{2010}{m}. A \qty{100}{km} más allá,
las edades son $450/10.6 = 42.5$ y $100/10.6 = \qty{9.4}{Myr}$, las profundidades
4781 y \qty{3575}{m}: un escalón de unos \qty{1210}{m}. A \qty{1000}{km} más allá,
las edades son 127.4 y \qty{94.3}{Myr}, las profundidades 6450 y \qty{5899}{m}:
\qty{550}{m}. (Estas edades superan \qty{80}{Myr}, donde la ley de la raíz cuadrada
se va demasiado hacia abajo; la ley de placa daría un escalón menor.)
\end{solution}

\begin{solution}{exo:b1:rifts-ridges-margins:10}
(a) Con áreas iguales para cada edad, la profundidad media es la media de $d(t)$ sobre
$[0, A]$:
\[
  \frac1A\int_0^A\bigl(d_r + C\sqrt t\bigr)\mathrm dt = d_r + \frac CA\cdot\frac23A^{3/2}
  = d_r + \frac23C\sqrt A.
\]
(b) Para $A = \qty{180}{Myr}$: $2500 + \frac23\times350\times13.416 =
\qty{5630}{m}$; para \qty{120}{Myr}: $2500 + \frac23\times350\times10.954 =
\qty{5056}{m}$, \qty{574}{m} menos profunda. (c) Sea $r$ la subida del nivel del mar y
el fondo cargado se hunda $s = r\rho_w/(\rho_m - \rho_w)$. El agua, sin cambio de volumen,
debe llenar la misma profundidad sobre la misma área, así que $r + s = \qty{574}{m}$,
$r\,\rho_m/(\rho_m - \rho_w) = 574$, y $r = 574\times2300/3300 =
\qty{400}{m}$: cientos de metros, solo por la velocidad de expansión.
\end{solution}

\begin{solution}{exo:b1:rifts-ridges-margins:11}
(a) $\sqrt{x - a} = \sqrt x\sqrt{1 - a/x} \approx \sqrt x\,(1 - a/2x)$ para
$a \ll x$, así que
$\Delta d \approx C\sqrt{x/u}\cdot a/(2x) = Ca/(2\sqrt{ux})$.
(b) En $x = \qty{450}{km}$: exacto \qty{1205}{m}, aproximado
$350\times350/(2\sqrt{10.6\times450}) = \qty{887}{m}$; en \qty{1350}{km}:
\qty{550}{m} y \qty{512}{m}; en \qty{3350}{km}: \qty{334}{m} y
\qty{325}{m}. El cociente entre el escalón aproximado y el exacto depende solo de
$r = a/x$: $(r/2)/(1 - \sqrt{1 - r})$, que vale 0.9 para $r = 0.36$. La
aproximación es buena al \qty{10}{\percent} más allá de $x = 350/0.36 \approx
\qty{970}{km}$, unos tres desplazamientos desde la dorsal.
\end{solution}

\begin{solution}{exo:b1:rifts-ridges-margins:12}
(a) $d_2 - d_1 = 3671.1 - 3606.8 = \qty{64.3}{m}$ a \qty{10}{Myr};
$4072.8 - 4065.3 = \qty{7.5}{m}$ a 20; $4415.4 - 4417.0 = \qty{-1.7}{m}$ a
30; $4707.5 - 4713.6 = \qty{-6.1}{m}$ a 40. (b) $d_2 - d_1$ es positiva a 20
y negativa a 30, y ambas leyes son continuas: son iguales entre ambas.
Bisección: a 25, $d_2 - d_1 = +0.9$; a 27.5, $-0.7$; a 26.25, $+0.01$; a
26.875, $-0.3$; a 26.56, $-0.2$; a 26.41, $-0.08$; a 26.33, $-0.04$: la
raíz yace entre 26.25 y 26.33, en \qty{26.3}{Myr}, donde ambas dan
\qty{4294}{m}. (c) $d_1' = 175/\sqrt{26.3} = \qty{34.1}{m.Myr^{-1}}$ y
$d_2' = (3200/62.8)\,e^{-26.3/62.8} = \qty{33.5}{m.Myr^{-1}}$. Las leyes tienen
casi el mismo valor y la misma pendiente allí, así que la curva empalmada no tiene
un quiebro visible.
\end{solution}

\begin{solution}{pb:b1:rifts-ridges-margins:1}
\textbf{1.} \omterm{def:b1:cross-sections:throw}{Salto horizontal} $= $ salto$/\tan$(\omterm{def:b1:geological-maps:attitude}{buzamiento}): F1 \qty{346}{m}, F2 \qty{770}{m},
F3 \qty{2014}{m}, F4 \qty{163}{m}, F5 \qty{462}{m}, F6 \qty{1172}{m}. F3, con
el mayor salto y \omterm{def:b1:cross-sections:throw}{salto horizontal}, es la falla de borde.

\textbf{2.} La extensión es la suma, \qty{4927}{m}, unos \qty{4.9}{km}. El
rift tenía $60 - 4.9 = \qty{55.1}{km}$ de anchura antes del fallamiento.

\textbf{3.} La normal a un rift orientado \ang{020} apunta a \ang{110}; el
movimiento, hacia \ang{105}, forma \ang{5} con ella. La componente perpendicular es
$6.0\cos\ang{5} = \qty{5.98}{mm.yr^{-1}}$.

\textbf{4.} $\qty{4.927}{km}/\qty{5.98}{km.Myr^{-1}} = \qty{0.82}{Myr}$.

\textbf{5.} En \qty{3.2}{Myr} las placas se separaron
$5.98\times3.2 = \qty{19.1}{km}$ a través del rift. Las fallas llevan
$4.93/19.1 = \qty{26}{\percent}$ de ello.

\textbf{6.} Los otros tres cuartos, \qty{14.2}{km}, no están en las fallas de la sección. Pueden absorberse por diques intruidos a lo
largo del eje, que abren el rift sin fallarlo, como en Afar, o por fallas fuera de la sección. Con diques de \qty{8}{m} de anchura,
$\qty{14200}{m}/\qty{8}{m} \approx 1780$ diques hacen falta, uno cada $3.2\times10^6/1780 \approx 1800$ años de media.

\textbf{7.} Desde C2An: $(21.0 + 20.2)/2.581 = 41.2/2.581 =
\qty{15.96}{mm.yr^{-1}}$. Desde C3n: $(41.0 + 40.2)/5.235 = 81.2/5.235 =
\qty{15.51}{mm.yr^{-1}}$.

\textbf{8.} $(41.2 - 2.2)/2.581 = \qty{15.11}{mm.yr^{-1}}$ y
$(81.2 - 2.2)/5.235 = \qty{15.09}{mm.yr^{-1}}$: corregidas, las dos concuerdan dentro de
\qty{0.02}{mm.yr^{-1}}. La corrección importa más para el cron más joven, cuya
distancia es menor.

\textbf{9.} El desplazamiento hacia fuera desplaza ambos bordes de un flanco en la misma cantidad, así que se cancela en su diferencia.
Sumando ambos flancos,
$(81.2 - 41.2)/(5.235 - 2.581) = 40.0/2.654 = \qty{15.07}{mm.yr^{-1}}$.

\textbf{10.} Nordeste: $(41.0 - 21.0)/2.654 = \qty{7.54}{mm.yr^{-1}}$;
suroeste: $(40.2 - 20.2)/2.654 = \qty{7.54}{mm.yr^{-1}}$. La expansión es
simétrica; la diferencia de \qty{0.8}{km} entre los flancos, la misma para los dos crones, es un desplazamiento constante, como si el eje
estuviera colocado \qty{0.4}{km} fuera del centro.

\textbf{11.} $\cos\Delta = \sin\ang{30.9}\sin\ang{18.0} +
\cos\ang{30.9}\cos\ang{18.0}\cos\ang{16.4} = 0.15870 + 0.78287 = 0.94157$, así que
$\Delta = \ang{19.69}$. Con $\omega R = 0.403\times(\pi/180)\times6371 =
\qty{44.81}{mm.yr^{-1}}$, la velocidad es $44.81\sin\ang{19.69} =
\qty{15.10}{mm.yr^{-1}}$, las velocidades corregidas de las preguntas 8 y 9.

\textbf{12.} $\cos\Delta = \sin\ang{30.9}\sin\ang{26.0} +
\cos\ang{30.9}\cos\ang{26.0}\cos\ang{11.4} = 0.22512 + 0.75600 = 0.98112$, así que
$\Delta = \ang{11.15}$ y la velocidad es $44.81\sin\ang{11.15} =
\qty{8.7}{mm.yr^{-1}}$.

\textbf{13.} Las predicciones, 15.1 y \qty{8.7}{mm.yr^{-1}}, concuerdan con las medidas 15 y \qty{9}{mm.yr^{-1}}. La velocidad es
$\omega R\sin\Delta$, y el polo yace al noroeste del mar Rojo, en el nordeste de Libia: cuanto más al sur, más lejos del polo, y más
rápida la apertura.

\textbf{14.} $15.09\times5.235 = \qty{79}{km}$ de fondo oceánico.

\textbf{15.} $5000/120.6 = \qty{41}{km.Myr^{-1}}$, \qty{41}{mm.yr^{-1}} de
media.

\textbf{16.} $5000/15.1 = \qty{331}{Myr}$ con la velocidad central;
$5000/8.7 = \qty{575}{Myr}$ con la velocidad norte.

\textbf{17.} Al abrirse alrededor de un polo cercano, el mar Rojo se ensancha como unas tijeras: su anchura crece en proporción a
$\sin\Delta$, y sigue siendo una cuña que se estrecha hacia el polo. A \qty{8.7}{mm.yr^{-1}}, la parte norte necesita
$79/8.7 = \qty{9.1}{Myr}$ para formar los \qty{79}{km} de fondo oceánico que la parte central tiene ahora.

\textbf{18.} A \qty{15}{mm.yr^{-1}} el eje del mar Rojo es una \omterm{def:b1:rifts-ridges-margins:ridge-classes}{dorsal ultralenta}; se espera un valle axial entre fallas normales.

\textbf{19.} $2500 + 350\times1 = \qty{2850}{m}$;
$2500 + 350\sqrt{10} = \qty{3607}{m}$; $2500 + 350\sqrt{40} = \qty{4714}{m}$.

\textbf{20.} $t = (1500/350)^2 = \qty{18.4}{Myr}$.

\textbf{21.} En los primeros \qty{10}{Myr} se hunde $350\sqrt{10} = \qty{1107}{m}$;
entre 70 y \qty{80}{Myr}, $350(\sqrt{80} - \sqrt{70}) = 350\times(8.944 -
8.367) = \qty{202}{m}$, cinco veces menos en el mismo lapso de tiempo.

\textbf{22.} $\mathrm dd/\mathrm dt = C/(2\sqrt t) = 175/\sqrt{80} =
\qty{19.6}{m.Myr^{-1}}$, unos \qty{0.02}{mm} al año.

\textbf{23.} $490/\sqrt{80} = \qty{54.8}{mW.m^{-2}}$.

\textbf{24.} Ley de la raíz cuadrada: $2500 + 350\sqrt{80} = 2500 + 350\times8.944 =
\qty{5630}{m}$. Ley de placa: $6400 - 3200\,e^{-80/62.8} = 6400 - 3200\times0.2798
= \qty{5505}{m}$. Difieren en \qty{126}{m}. A \qty{80}{Myr} ambas leyes siguen
aplicándose: la ley de la raíz cuadrada hasta unos \qty{80}{Myr}, la ley de placa desde
unos \qty{20}{Myr}; más allá, la ley de la raíz cuadrada se va demasiado hacia abajo.

\textbf{25.} La ley de la raíz cuadrada da \qty{5630}{m}: con precisión de una décima de kilómetro, el fondo del océano nuevo yace a una
profundidad de \textbf{\qty[text-series-to-math=true]{5.6}{km}} cuando tiene 80 millones de años (la ley de placa da \qty{5.5}{km}).
\end{solution}
